\section*{Introducción a la estadística inferencial}

La \emph{estadística inferencial} es la rama de la estadística que se ocupa de obtener conclusiones sobre una población a partir del análisis de una muestra. A diferencia de la estadística descriptiva, que se limita a resumir y describir los datos observados, la inferencia estadística nos permite generalizar y hacer predicciones más allá de los datos disponibles.

\subsection*{¿Por qué es necesaria la inferencia estadística?}

En la práctica, rara vez podemos estudiar a toda la población de interés debido a:

\begin{itemize}
	\item \textbf{Costo:} Estudiar a todos los individuos puede ser prohibitivamente caro.
	\item \textbf{Tiempo:} Recopilar datos de toda la población puede tomar demasiado tiempo.
	\item \textbf{Accesibilidad:} Algunos miembros de la población pueden ser inaccesibles.
	\item \textbf{Destrucción:} En algunos casos, la medición destruye el objeto estudiado (por ejemplo, probar la vida útil de una bombilla).
\end{itemize}

Por estas razones, tomamos una \emph{muestra representativa} y utilizamos técnicas estadísticas para inferir características de la población completa.

\subsection*{Conceptos fundamentales}

Antes de profundizar en las técnicas de inferencia, es crucial entender estos conceptos:

\begin{definicion}[Población y muestra]
	\begin{itemize}
		\item \textbf{Población:} Conjunto completo de individuos, objetos o eventos que poseen al menos una característica común que deseamos estudiar.
		\item \textbf{Muestra:} Subconjunto de la población seleccionado para el estudio.
	\end{itemize}
\end{definicion}

\begin{definicion}[Parámetro y estadístico]
	\begin{itemize}
		\item \textbf{Parámetro:} Medida numérica que describe una característica de la \emph{población}. Por ejemplo, la media poblacional $\mu$ o la desviación estándar poblacional $\sigma$.
		\item \textbf{Estadístico:} Medida numérica calculada a partir de una \emph{muestra}. Por ejemplo, la media muestral $\bar{x}$ o la desviación estándar muestral $s$.
	\end{itemize}
\end{definicion}

\begin{observacion}
	Los parámetros son valores fijos pero generalmente desconocidos. Los estadísticos son variables aleatorias que varían de muestra en muestra y se utilizan para estimar los parámetros.
\end{observacion}

\subsection*{Tipos de inferencia estadística}

Existen dos tipos principales de inferencia estadística:

\begin{enumerate}
	\item \textbf{Estimación:} Consiste en utilizar estadísticos muestrales para estimar el valor de un parámetro poblacional. Puede ser:
	\begin{itemize}
		\item \emph{Estimación puntual:} Proporcionar un único valor como estimación del parámetro.
		\item \emph{Estimación por intervalo:} Proporcionar un rango de valores (intervalo de confianza) dentro del cual esperamos que se encuentre el parámetro con cierto nivel de confianza.
	\end{itemize}
	
	\item \textbf{Pruebas de hipótesis:} Consiste en formular una afirmación (hipótesis) sobre un parámetro poblacional y utilizar datos muestrales para decidir si la evidencia respalda o contradice dicha afirmación.
\end{enumerate}

\subsection*{El proceso de inferencia estadística}

El proceso general de inferencia estadística sigue estos pasos:

\begin{enumerate}
	\item \textbf{Definir la población} de interés y la característica a estudiar.
	\item \textbf{Seleccionar una muestra representativa} mediante técnicas de muestreo apropiadas.
	\item \textbf{Calcular estadísticos muestrales} que servirán como base para la inferencia.
	\item \textbf{Aplicar técnicas de inferencia} (estimación o pruebas de hipótesis) utilizando distribuciones de probabilidad apropiadas.
	\item \textbf{Interpretar los resultados} en el contexto del problema original, considerando la incertidumbre inherente al proceso de muestreo.
\end{enumerate}

\subsection*{Fuentes de error en la inferencia}

Al hacer inferencias a partir de muestras, debemos considerar dos tipos de error:

\begin{itemize}
	\item \textbf{Error de muestreo:} Diferencia natural entre el estadístico muestral y el parámetro poblacional debido a la variabilidad aleatoria inherente al proceso de muestreo. Este error disminuye al aumentar el tamaño de la muestra.
	
	\item \textbf{Error de no muestreo:} Errores sistemáticos que no se deben al muestreo, como sesgo de selección, errores de medición, no respuesta, etc. Estos errores no disminuyen al aumentar el tamaño de la muestra y pueden ser más graves que los errores de muestreo.
\end{itemize}

\begin{observacion}
	La estadística inferencial se ocupa principalmente de cuantificar y controlar el error de muestreo. El control de errores de no muestreo requiere un diseño cuidadoso del estudio y procedimientos de recolección de datos apropiados.
\end{observacion}

\subsection*{Distribuciones muestrales}

Un concepto fundamental en inferencia estadística es el de \emph{distribución muestral}. Si tomamos muchas muestras del mismo tamaño de una población y calculamos un estadístico (como la media) para cada muestra, la distribución de estos estadísticos se llama distribución muestral.

El \textbf{Teorema del Límite Central} establece que, bajo ciertas condiciones, la distribución muestral de la media se aproxima a una distribución normal, independientemente de la forma de la distribución poblacional original, siempre que el tamaño de muestra sea suficientemente grande.

Este teorema es la piedra angular de muchas técnicas de inferencia estadística que estudiaremos en este capítulo.

\subsection*{Estructura del capítulo}

En las siguientes secciones estudiaremos:

\begin{enumerate}
	\item \textbf{Muestreo aleatorio y el Teorema del Límite Central:} Fundamentos teóricos para la inferencia.
	\item \textbf{Pruebas de hipótesis:} Marco conceptual y metodología para tomar decisiones basadas en datos.
	\item \textbf{Estadísticos Z y t:} Herramientas específicas para pruebas de hipótesis sobre medias.
	\item \textbf{Intervalos de confianza:} Método para estimar parámetros con un nivel de confianza cuantificado.
	\item \textbf{Prueba chi-cuadrada:} Técnica para analizar datos categóricos y probar independencia.
	\item \textbf{Correlación:} Medida de la relación lineal entre variables.
\end{enumerate}

A lo largo del capítulo, utilizaremos ejemplos prácticos y código en Python para ilustrar la aplicación de estos conceptos en situaciones reales.
